\subsection{Spatial populations and a common scalar continuum action}
\label{sec:source-metric-scalar-continuum}

The completed conservative source-record carrier is a Euclidean
three-space under the stated response and metric hypotheses. Its whole
dense record group is not a locally finite population. A finite-resolution
selection and a transport law can be stated on that carrier without
inserting a Cartesian lattice. This construction keeps both choices explicit.

Fix $r,\delta,\eta>0$, a response-metric ball $\Omega=\overline B(0,r)$,
and a finite directed
seam-word endpoint menu $M\subset\Omega$ that is an $\eta$-net of $\Omega$.
Such menus exist: density approximates a slightly inward-shifted point by
a source endpoint, compactness extracts finitely many approximants, and a
finite maximum word budget contains all of them. This is algebraic endpoint
availability, not a proof that a repair law accepts those words as moves.
Choose a maximal $\delta$-separated subset $S\subset M$, and set
$h=\delta+\eta$. The resolution, window and tie rule are supplied; maximal
means inclusion-maximal, not maximum cardinality or a selected energy minimum.

\begin{proposition}[Finite spatial population and metric transport]
\label{prop:source-metric-population}
The selected points form an $h$-net, and
\[
 (r/h)^3\le |S|\le(1+2r/\delta)^3.
\]
For $R\ge h$ and $B(x,R)\subset\Omega$,
\[
 ((R-h)/h)^3\le |S\cap B(x,R)|\le(1+2R/\delta)^3.
\]
Connect distinct points whose source distance is at most $a>2h$, weighting
edges by that distance. The graph is connected, degrees are bounded by
$(1+2a/\delta)^3-1$, and its shortest-path metric satisfies
\[
 |s-t|\le d_G(s,t)\le
 \left(1+\frac{2h}{a-2h}\right)|s-t|+2h.
\]
Consequently $a\to0$ and $h/a\to0$ give uniform metric convergence on the
fixed window. With $\eta\le\delta$ and $a=8\delta$, degrees instead stay
at most $4912$ and $d_G\le2|s-t|+4\delta$; this latter estimate alone
does not imply a Euclidean metric limit.
\end{proposition}
\begin{proof}
Maximality puts each menu point within $\delta$ of $S$, proving coverage.
Disjoint radius-$\delta/2$ balls give the upper counts; covering the window,
or $B(x,R-h)$, by radius-$h$ balls gives the lower counts. For the graph,
divide the segment from $s$ to $t$ into
$k=\lceil |s-t|/(a-2h)\rceil$ pieces. Convexity retains the segment in
$\Omega$; approximate its interior division points within $h$, keeping its
endpoints exact. Consecutive approximants are at most $a$ apart and the
path length is at most $|s-t|+2hk$. The triangle inequality gives the lower
bound. These estimates are independent of the choice among maximal sets.
\end{proof}

The metric fixes relative three-volume. Let $C_i$ be the Voronoi cell of
$x_i\in S$, clipped to $\Omega$, and $v_i=\operatorname{vol}(C_i)>0$.
Thus $\sum_i v_i=V=\operatorname{vol}(\Omega)$, and almost every point of
$C_i$ is at distance at most $h$ from $x_i$. These are metric weights,
not a physical event-count density. For $0<h\le\epsilon$ supply
\[
 a_{ij}=\frac{45v_iv_j}{\pi\epsilon^5}
       (1-|x_i-x_j|/\epsilon)_+,
 \qquad (Lq)_i=v_i^{-1}\sum_j a_{ij}(q_i-q_j).
\]
The weighted inner product is $\langle p,q\rangle_v=\sum_i v_ip_iq_i$.
Symmetry gives
$\langle q,Lq\rangle_v=\frac12\sum_{ij}a_{ij}(q_i-q_j)^2\ge0$.
The common finite action is
\begin{equation}
 S_h[q]=\frac12\int
 \left[\|\dot q\|_v^2-c^2\langle q,Lq\rangle_v-m^2\|q\|_v^2\right]dt,
 \qquad c>0,\quad m\ge0.
 \label{eq:source-metric-scalar-action}
\end{equation}
Its equation is $\ddot q+c^2Lq+m^2q=0$. The coefficient $m$ has inverse-time
units in this convention; neither physical mass nor time units are inferred.

\begin{proposition}[Uniform scalar and detector limit]
\label{prop:source-metric-scalar-limit}
Let $u_{tt}-c^2\Delta u+m^2u=0$ be a smooth reference on the full carrier.
On $[0,T]$, its spatial support has distance greater than $\epsilon$ from
the complement of $\Omega$. Write
$M_1=\sup|\nabla u|$ and $M_4=\sup\|D_x^4u\|_{\rm op}$, with suprema
over the full carrier and time window. Then
\[
 \big|(Lu|_S)_i+\Delta u(x_i)\big|
 \le \frac{15}{112}\epsilon^2M_4+
     \frac{120h(\epsilon+h)^3}{\epsilon^5}M_1
 \le \frac{15}{112}\epsilon^2M_4+\frac{960h}{\epsilon^2}M_1.
\]
Put $B=c^2\sqrt V[(15/112)\epsilon^2M_4+960h\epsilon^{-2}M_1]$.
Initialize the finite action with the sampled reference position and velocity.
Then
\[
 \|\dot q-u_t|_S\|_v\le tB,\qquad
 \|q-u|_S\|_v\le\tfrac12t^2B.
\]
For a fixed $C^1$ detector $b$, the field readout
$D_h=\sum_i v_ib(x_i)q_i$ and its continuum counterpart
$D=\int_\Omega bu$ obey
\[
 |D_h(t)-D(t)|\le
 \tfrac12\sqrt V\|b\|_\infty t^2B+
 Vh\|\nabla(bu(t))\|_\infty.
\]
Thus $\epsilon\to0$ and $h/\epsilon^2\to0$ yield a uniform finite-time
observable limit. In normalized carrier units $\epsilon=h^{1/4}$ gives
an $O(\sqrt h)$ bound. Nested populations are unnecessary.
\end{proposition}
\begin{proof}
Compare $L$ with the integral operator having kernel
$45(1-|z|/\epsilon)_+/(\pi\epsilon^5)$. Odd Taylor terms cancel and
\[
 \int (1-|z|/\epsilon)_+z_1^2\,dz=\frac{2\pi}{45}\epsilon^5,
 \qquad
 \int (1-|z|/\epsilon)_+|z|^4\,dz=\frac{\pi}{14}\epsilon^7.
\]
The second moment gives $-\Delta$; the fourth-order Taylor remainder,
bounded by $M_4|z|^4/24$, gives $15\epsilon^2M_4/112$.
For nearest-cell projection $P$, compare
$g_x(y)=(1-|x-y|/\epsilon)_+(u(x)-u(y))$ with $g_x(P(y))$.
Splitting the difference at the larger kernel value uses a point within
$\epsilon$ of $x$; kernel and field Lipschitz bounds give
$|g_x(P(y))-g_x(y)|\le2hM_1$. This difference vanishes outside
$B(x,\epsilon+h)$, giving the stated quadrature term. If $u(x)\ne0$,
the support buffer retains the whole kernel ball in $\Omega$; otherwise
the integrand outside $\Omega$ is zero. No boundary term is omitted.

The sampled residual in the weighted norm is at most $B$. The positive
operator $A=c^2L+m^2I$ has
$\|\cos(t\sqrt A)\|\le1$ and
$\|\sin(t\sqrt A)/\sqrt A\|\le t$, with the zero eigenspace interpreted
by continuity. Duhamel integration gives the two trajectory bounds.
Cauchy--Schwarz and nearest-cell quadrature for $bu$ give the detector bound.
\end{proof}

The same cell argument gives $\|L\|_v\le960/\epsilon^2$, using twice the
largest weighted degree. \texttt{MetricKernelEnergy} formalizes the finite
positive-energy and mass-weighted degree inequalities; the geometric,
integral and continuum arguments above are analytic proofs. A local
discrete-action time integrator with a strict Courant bound may consume
this operator and add its own time and initialization errors.

Smooth compact initial data admit the reference until their speed-$c$
domain of influence reaches the window boundary. Local energy flux for the
Klein--Gordon equation gives that speed; the nonnegative mass term changes
energy density but not the characteristic cone. Applied to a smooth local
preparation difference, the detector estimate bounds the finite response
outside that cone. A nonzero continuum signal is resolvable when its lower
bound exceeds the approximation and detector errors. Instantaneous small
tails of the finite graph are not identified with detectable superluminal
signals. This is an interior effective limit, not a theorem about arbitrary
boundary conditions or exact microscopic order reflection.

The canonical variables $Q_i=\sqrt{v_i}q_i$ and
$P_i=p_i/\sqrt{v_i}$, $p_i=v_i\dot q_i$, give the oscillator Hamiltonian
$\frac12(|P|^2+Q^T\widetilde A Q)$, where
$\widetilde A=M^{1/2}(c^2L+m^2I)M^{-1/2}\ge0$ and
$M=\operatorname{diag}(v_i)$. Supplied canonical quantization on
$L^2(\mathbb R^{|S|},dQ)$ gives the same field readouts as multiplication
operators. For $m>0$ every mode has a Gaussian vacuum; at $m=0$ zero
modes remain free particles unless a separate restriction is declared.
This finite Hilbert construction is not an interacting quantum continuum
or a derivation of the Standard Model.

Coarsening, the volume prescription, radial force law, inertia, action
coefficients and accessible observer variables remain added assumptions.
The kernel graph may have growing degree. Algebraic source words alone
do not implement its couplings as accepted primitive repairs or calibrate
their duration. The construction supplies one conditional scalar effective
world on the response carrier; its selection and physical interpretation
remain separate obligations.

\begin{proposition}[Controlled boost comparison for scalar detectors]
\label{prop:source-metric-boost-detectors}
Let $\Lambda$ be a proper orthochronous Lorentz transformation of
$c^2dt^2-|dx|^2$, and let $u$ be a smooth Klein--Gordon solution on the
full spacetime. Put $u^\Lambda=u\circ\Lambda^{-1}$. Suppose both
references satisfy the support and derivative hypotheses of
Proposition~\ref{prop:source-metric-scalar-limit} on the same $[0,T]$
and $\Omega$, with a common positive interior buffer. Choose a real
$f\in C_c^1(\mathbb R^{1+3})$ such that both $f$ and
$f^\Lambda=f\circ\Lambda^{-1}$ are supported in
$(0,T)\times\Omega^\circ$. For $w=u,u^\Lambda$, let $q_h^w$ solve
the same action~\eqref{eq:source-metric-scalar-action}, with the same
$L,v_i,c,m$, initialized by $w(0,x_i)$ and $\partial_tw(0,x_i)$.
Define the spacetime-smeared readouts
\[
 \mathcal D_h(f;q)=\int_0^T\sum_i v_i f(t,x_i)q_i(t)\,dt,
 \qquad
 \mathcal D(f;w)=\int_{\mathbb R^{1+3}}fw\,dt\,d^3x.
\]
Using the preceding spatial derivative bounds $M_1(w),M_4(w)$, put
\[
 B_w=c^2\sqrt V\left[
   \frac{15}{112}\epsilon^2M_4(w)
   +960h\epsilon^{-2}M_1(w)\right],
\]
and
\[
 \mathcal E_h(f,w)=
 \frac{\sqrt V B_w}{2}\int_0^T t^2\|f(t,\cdot)\|_\infty\,dt
 +Vh\int_0^T
       \|\nabla_x(f(t,\cdot)w(t,\cdot))\|_\infty\,dt.
\]
Then
\begin{equation}
 \left|\mathcal D_h(f^\Lambda;q_h^{u^\Lambda})
          -\mathcal D_h(f;q_h^u)\right|
 \le \mathcal E_h(f^\Lambda,u^\Lambda)+\mathcal E_h(f,u).
 \label{eq:source-metric-boost-detector-error}
\end{equation}
The right-hand side tends to zero for $\epsilon\to0$ and
$h/\epsilon^2\to0$. The conclusion is uniform over a bounded boost
class when the common window, buffer and displayed derivative and detector
bounds are uniform over that class.
\end{proposition}
\begin{proof}
Lorentz covariance of the Klein--Gordon equation makes $u^\Lambda$ a
solution with the same $c,m$. The preceding detector estimate applies at
each time with $b=f(t,\cdot)$; its time integral is
$|\mathcal D_h(f;q_h^w)-\mathcal D(f;w)|\le\mathcal E_h(f,w)$.
Since $\det\Lambda=1$ and both detector supports lie in the comparison
window, change of variables gives
$\mathcal D(f^\Lambda;u^\Lambda)=\mathcal D(f;u)$.
The two approximation bounds and the triangle inequality prove the claim.
\end{proof}

The transformed preparations are supplied through the continuum reference
on the transformed Cauchy surface; this comparison does not construct a
boost map from one finite history or its initial data. The same estimate
applies to smooth local preparation differences by linearity. Spacetime
smearing is essential to the displayed comparison: a fixed-time detector
instead transforms to a tilted hypersurface with its transformed measure.
The result controls detector values for one common declared action, not
exact finite-graph support, native source selection, or boosts of arbitrary
finite records. Time integration and preparation errors, if present in an
execution, require their own bounds.
